\begingroup\color{red}
\chapter{Analyse de la prestation}
Ce chapitre sera finalisé après l’action. Je préfère ne pas écrire à l’avance des points forts ou des difficultés qui n’ont pas encore été observés.

\section{Ce qui a fonctionné et ce qui a moins bien fonctionné}

Après la séance, je m’appuierai sur des faits précis : participation du groupe, compréhension des consignes, gestion du temps, qualité des échanges, fonctionnement du binôme et sécurité de l’activité extérieure.

[À compléter : Présenter deux ou trois faits précis observés le 4 décembre. Pour chacun : ce qui s’est passé, pourquoi cela a pu se produire, quel effet cela a eu sur la séance et ce que j’en retiens.]

\section{Retours des stagiaires et du commanditaire}

Je comparerai les résultats du quiz, les réponses au questionnaire de satisfaction et le retour de Mme LIÉGEOIS avec les attentes identifiées au début du projet.

[À compléter : Ajouter le nombre de réponses recueillies, les tendances principales, quelques remarques anonymisées si elles sont utiles, et le retour du commanditaire.]

\section{Ce que je changerais}

L’objectif de cette partie sera de montrer ce que j’ai appris comme futur ECSR. Les améliorations devront être directement liées aux constats : simplifier une consigne, modifier une durée, mieux intégrer les non-conducteurs, ajuster le support ou approfondir un débriefing.

[À compléter : Rédiger des améliorations concrètes et réalistes : ce que je garderais, ce que je modifierais et ce que je ferais différemment lors d’une prochaine action.]
\endgroup
